\documentclass[12pt,a4paper]{article}

%─── Packages ──────────────────────────────────────────────────────────────────
\usepackage[margin=2.5cm]{geometry}
\usepackage{amsmath,amssymb}
\usepackage{graphicx}
\usepackage{booktabs}
\usepackage{caption}
\usepackage{subcaption}
\usepackage{hyperref}
\usepackage{siunitx}
\usepackage{physics}
\usepackage{microtype}
\usepackage{parskip}
\usepackage[numbers,sort&compress]{natbib}
\usepackage{xcolor}
\usepackage{fancyhdr}
\usepackage{titlesec}

%─── Formatting ────────────────────────────────────────────────────────────────
\hypersetup{colorlinks=true, linkcolor=blue!60!black, citecolor=green!40!black,
            urlcolor=blue!60!black}
\sisetup{separate-uncertainty=true}
\captionsetup{font=small, labelfont=bf, width=0.92\linewidth}

\pagestyle{fancy}
\fancyhf{}
\lhead{\small UV-Vis Band Gap Estimation}
\rhead{\small \today}
\cfoot{\thepage}
\renewcommand{\headrulewidth}{0.4pt}

\titleformat{\section}{\large\bfseries}{}{0em}{}[\titlerule]
\titleformat{\subsection}{\normalsize\bfseries}{}{0em}{}

%─── Macros ────────────────────────────────────────────────────────────────────
\newcommand{\Eg}{E_{\mathrm{g}}}
\newcommand{\hnu}{h\nu}
\newcommand{\cm}{\,\mathrm{cm}^{-1}}

%───────────────────────────────────────────────────────────────────────────────
\begin{document}

%─── Title ─────────────────────────────────────────────────────────────────────
\begin{center}
  {\LARGE\bfseries UV-Vis Spectroscopy for Optical Band Gap Estimation}\\[0.6em]
  {\large Tauc Plot Analysis of Direct and Indirect Semiconductors}\\[1.2em]
  {\normalsize\itshape Laboratory Report}\\[0.3em]
  {\small \today}
\end{center}

\vspace{0.5em}
\noindent\rule{\linewidth}{0.6pt}
\vspace{0.4em}

%─── Abstract ──────────────────────────────────────────────────────────────────
\begin{abstract}
\noindent
The optical band gaps of four direct-gap semiconductor thin films (CdS, TiO$_2$, ZnO,
ZnTe) and one indirect-gap Si-based polymer film were estimated using the Tauc plot
method from UV-Vis absorbance data. The absorption coefficient $\alpha$ was extracted
via the Beer–Lambert law. For direct transitions ($\alpha^2 \propto \hnu - \Eg$), linear
extrapolation of $\alpha^2$ vs.\ $\hnu$ yields the band gap; for indirect transitions
($\sqrt{\alpha} \propto \hnu - \Eg$), the same procedure is applied to $\sqrt{\alpha}$.
Extracted band gaps: CdS $\approx \SI{1.81}{eV}$, ZnTe $\approx \SI{1.92}{eV}$,
ZnO $\approx \SI{3.61}{eV}$, and Si-polymer
$\Eg = \SI[parse-numbers=false]{1.345 \pm 0.071}{eV}$.
TiO$_2$ data were saturated at the detector limit.
Full error propagation and comparison with literature values are reported.
The Python code for the analysis is available at:
\begin{center}
  \url{https://github.com/YOUR_USERNAME/bandgap-uvvis}   % \textleftarrow{} fill in your repo URL
\end{center}
\end{abstract}

\vspace{0.4em}
\noindent\rule{\linewidth}{0.4pt}
\vspace{0.2em}

%───────────────────────────────────────────────────────────────────────────────
\section{Introduction}
%───────────────────────────────────────────────────────────────────────────────

The band gap $\Eg$ is the fundamental energy difference between the valence band
maximum (VBM) and conduction band minimum (CBM).  In a \textbf{direct} band gap
semiconductor the VBM and CBM occur at the same crystal momentum $\mathbf{k}$, so an
electron-hole pair can be created by a single photon of energy $\hnu \geq \Eg$.  In an
\textbf{indirect} gap material the extrema occur at different $\mathbf{k}$ values; absorption
therefore requires a simultaneous phonon interaction to conserve momentum, leading to a
weaker, gradual absorption onset~\cite{Pankove1971}.

UV-Vis spectroscopy probes electronic transitions as a function of photon energy and is a
standard, non-destructive technique for optical band gap determination in thin films.  The
Tauc plot method~\cite{Tauc1966} linearises the absorption edge to give a precise
x-intercept estimate of $\Eg$.

%───────────────────────────────────────────────────────────────────────────────
\section{Theory}
%───────────────────────────────────────────────────────────────────────────────

\subsection{Beer–Lambert Law and Absorption Coefficient}

For a thin film of thickness $d$, the Beer–Lambert law gives
\begin{equation}
  I_T = I_0 \, e^{-\alpha d},
  \label{eq:BeerLambert}
\end{equation}
where $I_0$ and $I_T$ are the incident and transmitted intensities.
The absorbance $A = \log_{10}(I_0/I_T)$, so the absorption coefficient is
\begin{equation}
  \alpha \;(\mathrm{cm}^{-1}) = \frac{2.303\,A}{d},
  \label{eq:alpha}
\end{equation}
with $d$ in cm.

\subsection{Tauc Relation}

Near the band edge the absorption coefficient follows~\cite{Tauc1966,Pankove1971}
\begin{equation}
  \alpha = C \bigl[\hnu - \Eg\bigr]^{n},
  \label{eq:Tauc}
\end{equation}
where $C$ is a material-dependent constant.  The exponent $n$ encodes the type of
transition:
\begin{itemize}
  \item $n = \tfrac{1}{2}$: direct allowed transitions $\;\Rightarrow\; \alpha^2 \propto (\hnu-\Eg)$
  \item $n = 2$: indirect allowed transitions $\;\Rightarrow\; \sqrt{\alpha} \propto (\hnu-\Eg)$
\end{itemize}
Plotting $\alpha^2$ (or $\sqrt{\alpha}$) versus $\hnu$ and linearly extrapolating the
absorption edge to the energy axis gives $\Eg$ as the x-intercept.

The square-root onset in the density of states $D_{cv}(E) \propto (E-\Eg)^{1/2}$ arises
from the 3-D parabolic band approximation: near the band edge $k^2 \propto (E-\Eg)$, and
the number of states inside a sphere of radius $k$ scales as $k \propto (E-\Eg)^{1/2}$.
The matrix element $M_{cv}(\mathbf{k}) \approx \text{const.}$ near $\mathbf{k}=\mathbf{0}$
(smooth function, Taylor-expanded to leading order), giving $\alpha(\hnu) \propto |M_{cv}|^2
D_{cv}(\hnu) \propto (\hnu - \Eg)^{1/2}$ for direct transitions.

%───────────────────────────────────────────────────────────────────────────────
\section{Experimental}
%───────────────────────────────────────────────────────────────────────────────

\subsection{Instrument}

A Holmarc fibre-optic spectrometer (CCD array, 3648 pixels, wavelength range
\SIrange{386.5}{995.6}{nm}) was used.  The exposure time was optimised to keep peak
counts below 60\,000 (saturation at 65\,535) with 10 averages.  Dark and reference
spectra were recorded before each sample.  Samples were inserted normal to the beam
to minimise reflection losses.

\subsection{Samples}

\begin{table}[h]
  \centering
  \caption{Samples analysed in this experiment.}
  \label{tab:samples}
  \begin{tabular}{llll}
    \toprule
    Sample & Type & $d$ (nm) & Gap type \\
    \midrule
    CdS    & thin film on glass & 200 & Direct \\
    TiO$_2$ & thin film on glass & 200 & Direct \\
    ZnO    & thin film on glass & 200 & Direct \\
    ZnTe   & thin film on glass & 200 & Direct \\
    Si-polymer & thin film, 150 nm & 150 & Indirect \\
    \bottomrule
  \end{tabular}
\end{table}

\subsection{Data Processing}

Pixel index was converted to wavelength via the linear calibration
$\lambda = 386.542 + (p-1)\times 0.1674$\,nm (instrument specification).
Photon energy: $E = hc/\lambda = 1239.8/\lambda$\,eV.
The absorption coefficient was computed via Eq.~\eqref{eq:alpha} using the
spectrometer-exported absorbance.  For the indirect sample, data below
$\SI{1.2}{eV}$ were excluded due to an impurity-related absorption shoulder
(analogous to Fig.~2 in the lab manual for InSb~\cite{labmanual}).
All analysis was performed with a custom Python 3 script; the code is available at
the GitHub repository listed in the abstract.

%───────────────────────────────────────────────────────────────────────────────
\section{Results}
%───────────────────────────────────────────────────────────────────────────────

\subsection{Absorbance Spectra}

Figure~\ref{fig:absorbance} shows the absorbance spectra of all direct-gap samples.
CdS exhibits substantial background absorption across the entire visible range,
consistent with a thick/defective film.  ZnTe shows a clear absorption onset near
$\sim\SI{590}{nm}$ ($\sim\SI{2.1}{eV}$).  ZnO and TiO$_2$ absorb strongly in the
UV, with band edges at or beyond the short-wavelength detector limit ($\sim\SI{387}{nm}$,
$\SI{3.21}{eV}$).

\begin{figure}[htbp]
  \centering
  \includegraphics[width=0.82\linewidth]{figures/fig1_absorbance.png}
  \caption{Absorbance vs.\ wavelength for the four direct-gap samples.  ZnO and
  TiO$_2$ band edges extend beyond the detector range ($>\SI{3.21}{eV}$).}
  \label{fig:absorbance}
\end{figure}

\subsection{Direct Band Gap — Tauc Plots ($\alpha^2$ vs.\ $h\nu$)}

Figure~\ref{fig:tauc_direct} shows the combined Tauc plots.  Individual plots with
annotated linear fits are given in Appendix~\ref{app:individual}.

\begin{figure}[htbp]
  \centering
  \includegraphics[width=0.92\linewidth]{figures/fig2_tauc_direct.png}
  \caption{Tauc plots ($\alpha^2$ vs.\ $\hnu$) for all direct-gap samples.
  Red dashed lines: linear fits in the identified onset region.
  TiO$_2$: Tauc plot shown but Eg not extracted (see text).}
  \label{fig:tauc_direct}
\end{figure}

\paragraph{CdS.}  The film absorbs throughout the measured range, indicating
significant subgap absorption (likely from defects/impurities or a thick film).
A linear region is identified at \SIrange{2.10}{2.40}{eV}, yielding
$\Eg(\mathrm{CdS}) \approx \SI{1.81}{eV}$ ($R^2 = 0.97$).  The literature value
is \SI{2.42}{eV}; the discrepancy is attributed to background absorption raising the
$\alpha^2$ baseline, which shifts the x-intercept to lower energies.

\paragraph{ZnTe.}  A clean linear onset is observed at \SIrange{2.10}{2.40}{eV},
giving $\Eg(\mathrm{ZnTe}) \approx \SI{1.92}{eV}$ ($R^2 = 0.97$).
Literature: \SI{2.26}{eV}.

\paragraph{ZnO.}  The band edge falls near the detector limit; only a few points
above \SI{3.0}{eV} define the linear region, yielding
$\Eg(\mathrm{ZnO}) \approx \SI{3.61}{eV}$.
Literature: \SI{3.37}{eV}; the agreement is reasonable given the limited spectral range.

\paragraph{TiO$_2$.}  The absorbance is constant ($A \approx 0.23$) above
$\sim\SI{2.9}{eV}$, indicating detector saturation in the UV, and the data are
strongly noise-dominated below that.  No reliable Tauc fit could be extracted; the
true band gap ($\SI{3.20}{eV}$) lies at the spectrometer's short-wavelength
boundary.

\subsection{Indirect Band Gap — Tauc Plot ($\sqrt{\alpha}$ vs.\ $h\nu$)}

Figure~\ref{fig:alpha_indirect} shows $\alpha$ vs.\ $\hnu$ for the Si-based film;
Figure~\ref{fig:tauc_indirect} shows the Tauc plot.

\begin{figure}[htbp]
  \centering
  \begin{subfigure}[b]{0.48\linewidth}
    \includegraphics[width=\linewidth]{figures/fig4_alpha_indirect.png}
    \caption{$\alpha$ vs.\ photon energy.  The gradual onset is characteristic of
    indirect transitions.}
    \label{fig:alpha_indirect}
  \end{subfigure}\hfill
  \begin{subfigure}[b]{0.48\linewidth}
    \includegraphics[width=\linewidth]{figures/fig3_tauc_indirect.png}
    \caption{Tauc plot ($\sqrt{\alpha}$ vs.\ $\hnu$).  Shaded region: excluded
    impurity shoulder ($<\SI{1.2}{eV}$).}
    \label{fig:tauc_indirect}
  \end{subfigure}
  \caption{Absorption data for the indirect-gap Si-based polymer film
  ($d = \SI{150}{nm}$).}
\end{figure}

A linear fit over \SIrange{1.2}{2.0}{eV} gives:
\begin{align}
  \sqrt{\alpha} &= 572.7\,(\hnu) - 770.3 \qquad (R^2 = 0.803) \\[4pt]
  \Eg &= \frac{770.3}{572.7} = \SI{1.345}{eV}
\end{align}
The $R^2 = 0.80$ reflects the inherent curvature in indirect-gap absorption (phonon
branches contribute two slightly different slopes; cf.\ Fig.~1 of the lab
manual~\cite{labmanual}).

%───────────────────────────────────────────────────────────────────────────────
\section{Error Analysis}
%───────────────────────────────────────────────────────────────────────────────

\subsection{Propagation through $\alpha$}

From Eq.~\eqref{eq:alpha}, applying partial-derivative propagation:
\begin{equation}
  \frac{\delta\alpha}{\alpha} = \sqrt{
    \left(\frac{\delta A}{A}\right)^2 + \left(\frac{\delta d}{d}\right)^2
  },
  \label{eq:dalpha}
\end{equation}
with $\delta A = 0.005$ (instrument resolution) and $\delta d = \SI{5}{nm}$
(deposition tolerance, $\sim 3.3\%$).  At the indirect-gap onset
($A \approx 0.3$, $d = \SI{150}{nm}$), both terms contribute equally,
giving $\delta\alpha/\alpha \approx 4.7\%$.

\subsection{Propagation through $\sqrt{\alpha}$}

By the chain rule: $\delta(\sqrt{\alpha}) = \delta\alpha / (2\sqrt{\alpha})$.

\subsection{Uncertainty in $\Eg$}

The band gap is extracted as $\Eg = -b/m$ from a linear fit $\sqrt{\alpha} = m\,\hnu + b$.
Error propagation through the ratio:
\begin{equation}
  \frac{\delta\Eg}{\Eg} = \sqrt{
    \left(\frac{\delta m}{m}\right)^2 + \left(\frac{\delta b}{b}\right)^2
  },
\end{equation}
where $\delta m$ and $\delta b$ are the ordinary least-squares standard errors of the
slope and intercept.  For the indirect sample:

\begin{table}[h]
  \centering
  \caption{Uncertainty budget for indirect band gap.}
  \label{tab:errors}
  \begin{tabular}{llS[table-format=2.4]}
    \toprule
    Source & Symbol & {Value} \\
    \midrule
    Absorbance resolution & $\delta A$ & 0.0050 \\
    Thickness uncertainty & $\delta d$ & \SI{5}{nm} \\
    Std.\ error of slope  & $\delta m$ & \SI{21.37}{(cm^{-1})^{1/2} eV^{-1}} \\
    Std.\ error of intercept & $\delta b$ & \SI{37.63}{(cm^{-1})^{1/2}} \\
    \midrule
    \textbf{Band gap uncertainty} & $\delta\Eg$ & \SI{0.071}{eV} \\
    \bottomrule
  \end{tabular}
\end{table}

\noindent
\textbf{Final result (indirect):}
\begin{equation}
  \boxed{\Eg = \SI[parse-numbers=false]{1.345 \pm 0.071}{eV}}
\end{equation}

\subsection{Dominant Error Sources}

\begin{enumerate}
  \item \textbf{Regression statistical error} ($\delta m$, $\delta b$) — dominant.
        Caused by curvature in the Tauc plot; a wider or better-defined linear region
        would reduce this.
  \item \textbf{Thickness uncertainty} — equally important with absorbance error
        at the band onset (Eq.~\eqref{eq:dalpha}).
  \item \textbf{Linear region selection} — systematic and largest for CdS/ZnTe,
        where background absorption displaces the intercept.
  \item \textbf{Detector saturation / spectral range} — limits TiO$_2$ and ZnO
        analysis entirely.
  \item \textbf{Reflection and scattering losses} — overestimate $\alpha$
        (not corrected here).
  \item \textbf{Parabolic band approximation} — model error inherent to the Tauc
        method; real bands deviate from perfect parabolicity.
\end{enumerate}

%───────────────────────────────────────────────────────────────────────────────
\section{Discussion}
%───────────────────────────────────────────────────────────────────────────────

\begin{table}[htbp]
  \centering
  \caption{Summary of measured band gaps vs.\ literature values at \SI{300}{K}.}
  \label{tab:summary}
  \begin{tabular}{lccc S[table-format=+3.1] l}
    \toprule
    Sample & Type & $\Eg^{\mathrm{meas}}$ (eV) & $\Eg^{\mathrm{lit}}$ (eV)
           & {$\Delta$ (\%)} & Remark \\
    \midrule
    CdS     & Direct   & 1.814 & 2.42\cite{Hummel2001} & -25.1 & background abs. \\
    TiO$_2$ & Direct   & N/A   & 3.20\cite{Hummel2001} & {—}   & detector limited \\
    ZnO     & Direct   & 3.614 & 3.37\cite{Hummel2001} & +7.2  & near limit \\
    ZnTe    & Direct   & 1.918 & 2.26\cite{Hummel2001} & -15.1 & background abs. \\
    Si-polymer & Indirect & $1.345\pm0.071$ & 1.12\cite{Hummel2001} & +20.1 & \\
    \bottomrule
  \end{tabular}
\end{table}

For CdS and ZnTe the measured $\Eg$ is lower than literature because the films absorb
even below the true band edge (defect states, impurities, or interference from the glass
substrate), which raises the $\alpha^2$ baseline and shifts the linear extrapolation
intercept to lower energies.  This is a known artefact of thick or defect-rich films.

The Si-based polymer gives $\Eg = \SI{1.345}{eV}$, somewhat above crystalline Si
(\SI{1.12}{eV}).  This is consistent with quantum confinement or composition effects in a
polymer matrix~\cite{labmanual}, or with the $R^2 = 0.80$ suggesting a non-ideal
linear region.

ZnO gives $+7\%$ deviation; because only a handful of data points lie near the band edge,
the fit is sensitive to noise, but the order-of-magnitude agreement confirms the measurement
is physically meaningful.

The gradual vs.\ sharp absorption onset in Fig.~\ref{fig:alpha_indirect} vs.\
Fig.~\ref{fig:absorbance} qualitatively confirms the distinction between indirect
(phonon-mediated, slow) and direct (fast, sharp) transitions.

%───────────────────────────────────────────────────────────────────────────────
\section{Conclusion}
%───────────────────────────────────────────────────────────────────────────────

The Tauc plot method was successfully applied to UV-Vis absorbance data to extract optical
band gaps.  For ZnTe and ZnO, where the absorption edge falls within the spectrometer
range, the results agree with literature to within 7--15\%.  The indirect band gap of the
Si-polymer film is $\Eg = \SI[parse-numbers=false]{1.345 \pm 0.071}{eV}$, in reasonable
agreement with Si-based materials.  CdS and TiO$_2$ are limited by background absorption
and detector range, respectively, illustrating the importance of matching the spectrometer
range to the material's band gap.

%───────────────────────────────────────────────────────────────────────────────
\begin{thebibliography}{9}

\bibitem{Pankove1971}
  J.~I.~Pankove,
  \textit{Optical Processes in Semiconductors},
  Dover Publications, 1971.

\bibitem{Tauc1966}
  J.~Tauc, R.~Grigorovici, A.~Vancu,
  \textit{Phys.\ Status Solidi}, \textbf{15}, 627 (1966).

\bibitem{Hummel2001}
  R.~E.~Hummel,
  \textit{Electronic Properties of Materials}, 3rd edn.,
  Springer, 2001, Appendix 4.

\bibitem{Zanatta2019}
  A.~R.~Zanatta,
  \textit{Sci.\ Rep.}, \textbf{9}, 11225 (2019).

\bibitem{labmanual}
  Lab Manual,
  \textit{Ultraviolet-Visible Spectroscopy to Determine the Nature of Band Gap
  and Estimation of the Bandgap},
  IIT, 2024.

\end{thebibliography}

%───────────────────────────────────────────────────────────────────────────────
\appendix
\section{Individual Tauc Plots}
\label{app:individual}
%───────────────────────────────────────────────────────────────────────────────

\begin{figure}[htbp]
  \centering
  \includegraphics[width=\linewidth]{figures/fig5_cds_tauc.png}
  \caption{CdS: absorbance spectrum (left) and Tauc plot with linear fit (right).}
\end{figure}

\begin{figure}[htbp]
  \centering
  \includegraphics[width=\linewidth]{figures/fig6_tio2_tauc.png}
  \caption{TiO$_2$: absorbance spectrum (left) and Tauc plot (right). Reliable
  band gap extraction is not possible due to data saturation.}
\end{figure}

\begin{figure}[htbp]
  \centering
  \includegraphics[width=\linewidth]{figures/fig7_zno_tauc.png}
  \caption{ZnO: absorbance spectrum (left) and Tauc plot with linear fit (right).}
\end{figure}

\begin{figure}[htbp]
  \centering
  \includegraphics[width=\linewidth]{figures/fig8_znte_tauc.png}
  \caption{ZnTe: absorbance spectrum (left) and Tauc plot with linear fit (right).}
\end{figure}

\end{document}
